% Add to the preamble of the main document:
%   \usepackage{tikz}
%   \usetikzlibrary{positioning, arrows.meta}
%
% Then insert the section below anywhere before the "Negative Results" section.

\section{The Toy Model}
\label{sec:toy-model}

We study a controlled sequence-classification task defined on a random layered
directed graph. The graph has $L+1$ node layers $V_0, V_1, \dots, V_L$, with
$|V_0|=1$ (a single seed node) and $|V_\ell| = N$ for $\ell\geq 1$. Each node in
layer $\ell < L$ has exactly $k$ outgoing edges, sampled without replacement
from $V_{\ell+1}$. The set of all paths from the seed to layer $L$ defines the
dataset. Each example is an edge sequence
\begin{equation*}
    x = (e_0, e_1, \dots, e_{L-1}) \in \{0, 1, \dots, k-1\}^{L},
\end{equation*}
where $e_\ell \in \{0, \dots, k-1\}$ selects one of the $k$ outgoing edges of
the node reached at layer $\ell$. The label is the identity of the node reached
in $V_L$:
\begin{equation*}
    y = z_L(x) \in V_L.
\end{equation*}

Along each path, the intermediate nodes $z_1(x), \dots, z_{L-1}(x)$ are
\emph{latent variables}: each $z_\ell$ is a deterministic function of the
prefix $(e_0, \dots, e_{\ell-1})$, but the model only observes the raw edge
sequence and must implicitly recover the $z_\ell$ in order to compose them
into the final label $z_L$.

\begin{figure}[h]
\centering
\begin{tikzpicture}[
    every node/.style={circle, draw, minimum size=6mm, inner sep=0pt,
                       font=\scriptsize},
    edge/.style={-{Latex[length=1.6mm]}, gray!55, thin},
    hi/.style={-{Latex[length=1.8mm]}, red!80!black, very thick},
    lbl/.style={draw=none, font=\itshape}
]
% --- Nodes: 5 layers, 1 -> 4 -> 4 -> 4 -> 4, k = 2 ---
\node (a0) at (0, 0)   {$s$};

\node (b0) at (2, 1.8) {$z_1^{(0)}$};
\node (b1) at (2, 0.6) {$z_1^{(1)}$};
\node (b2) at (2,-0.6) {$z_1^{(2)}$};
\node (b3) at (2,-1.8) {$z_1^{(3)}$};

\node (c0) at (4, 1.8) {};
\node (c1) at (4, 0.6) {};
\node (c2) at (4,-0.6) {};
\node (c3) at (4,-1.8) {};

\node (d0) at (6, 1.8) {};
\node (d1) at (6, 0.6) {};
\node (d2) at (6,-0.6) {};
\node (d3) at (6,-1.8) {};

\node (e0) at (8, 1.8) {};
\node (e1) at (8, 0.6) {};
\node (e2) at (8,-0.6) {$y$};
\node (e3) at (8,-1.8) {};

% --- Random edges (k = 2 per node, chosen for a clean illustration) ---
\draw[edge] (a0) -- (b0);  \draw[edge] (a0) -- (b2);
\draw[edge] (b0) -- (c1);  \draw[edge] (b0) -- (c2);
\draw[edge] (b1) -- (c0);  \draw[edge] (b1) -- (c3);
\draw[edge] (b2) -- (c1);  \draw[edge] (b2) -- (c3);
\draw[edge] (b3) -- (c0);  \draw[edge] (b3) -- (c2);
\draw[edge] (c0) -- (d0);  \draw[edge] (c0) -- (d1);
\draw[edge] (c1) -- (d2);  \draw[edge] (c1) -- (d3);
\draw[edge] (c2) -- (d0);  \draw[edge] (c2) -- (d3);
\draw[edge] (c3) -- (d1);  \draw[edge] (c3) -- (d2);
\draw[edge] (d0) -- (e0);  \draw[edge] (d0) -- (e2);
\draw[edge] (d1) -- (e1);  \draw[edge] (d1) -- (e3);
\draw[edge] (d2) -- (e0);  \draw[edge] (d2) -- (e1);
\draw[edge] (d3) -- (e2);  \draw[edge] (d3) -- (e3);

% --- One highlighted example path ---
\draw[hi] (a0) -- node[above, sloped, draw=none, font=\tiny]{$e_0$} (b0);
\draw[hi] (b0) -- node[above, sloped, draw=none, font=\tiny]{$e_1$} (c2);
\draw[hi] (c2) -- node[above, sloped, draw=none, font=\tiny]{$e_2$} (d3);
\draw[hi] (d3) -- node[above, sloped, draw=none, font=\tiny]{$e_3$} (e2);

% --- Layer labels ---
\node[lbl] at (0,-2.8) {$V_0$};
\node[lbl] at (2,-2.8) {$V_1$};
\node[lbl] at (4,-2.8) {$V_2$};
\node[lbl] at (6,-2.8) {$V_3$};
\node[lbl] at (8,-2.8) {$V_L$};
\end{tikzpicture}
\caption{A small instance with $L=4$ layers and $k=2$ edges per node. Each node
in $V_\ell$ has $k$ outgoing edges to $V_{\ell+1}$. The model observes the
edge sequence $(e_0, e_1, e_2, e_3)$ along a path (red) and predicts the final
node $y\in V_L$. The intermediate nodes visited in $V_1, V_2, V_3$ are the
latent variables $z_1, z_2, z_3$ whose emergence we probe.}
\label{fig:toy-graph}
\end{figure}

\subsection*{Causal geometry of the latents}

A path is fully specified by its edge sequence $x$, and the label $y$ is a
composition of $L$ successive edge lookups. This causal structure has two
immediate consequences for probing:

\begin{itemize}
    \item \textbf{Position.} The latent $z_\ell$ depends only on
    $(e_0, \dots, e_{\ell-1})$. In a causal transformer, information about
    $z_\ell$ can therefore first appear at input position $p = \ell - 1$ and
    only propagate to later positions through attention.

    \item \textbf{Depth.} Since $z_\ell$ composes $\ell$ successive edges, it
    typically requires at least $\lceil \log_2 \ell \rceil$ blocks of nonlinear
    composition to be resolved from the raw edge embeddings. In practice we
    observe that deeper graph latents emerge at deeper transformer depths.
\end{itemize}

\subsection*{Model and probing objective}

We train a small transformer $h_\theta : \{0,\dots,k-1\}^L \to \Delta(V_L)$ to
predict $y$ from $x$. The residual stream at transformer depth $d$ and
sequence position $p$ produces a vector $h_i^{(d,p)}(\theta)\in\mathbb{R}^m$
for each example $i$. This yields a grid of $(D+1)\times L$ candidate probe
sites, where $D$ is the number of transformer blocks.

Given a target graph layer $\ell$, a probe on cell $(d, p)$ is a classifier
$f: \mathbb{R}^m \to V_\ell$ trained on
$\{(h_i^{(d, p)}, z_\ell(x_i))\}_{i \in \text{train}}$ and evaluated on a
held-out split. We compare three families of probe:
\begin{itemize}
    \item \emph{Difference-of-means} (DoM), the rank-$1$ direction
    $w_c = \mu_c^+ - \mu_c^-$;
    \item \emph{Multinomial logistic regression}, a full-rank linear classifier;
    \item \emph{One-hidden-layer MLP}, a nonlinear classifier.
\end{itemize}

Comparing these three probes across training reveals when a latent becomes
decodable and whether it is first acquired in a linear or a nonlinear form.
The gap between DoM and logistic regression isolates whether the linearly
decodable content sits along a single mean-difference direction or is spread
across the full covariance structure. The gap between logistic regression and
MLP isolates whether an additional nonlinear coordinate change is required.

\subsection*{Default configuration}

Unless stated otherwise, the experiments use $L = 6$, $|V_\ell| = 100$ for
$\ell \geq 1$, and $k = 10$, giving $k^L = 10^6$ training paths. The
transformer has $4$ blocks with a learned positional embedding, causal
self-attention, and ReLU$^2$ MLPs; the classifier reads the residual at the
last sequence position. Probes are refit from scratch on the current model
every $50$ training steps and evaluated on a fixed test split, so a probe
curve $\{\mathrm{acc}(d, p; t)\}$ can be plotted against training step $t$.
